\section{Conclusion}
\label{conclusions}

In this paper, we investigated the effectiveness of the general-purpose CGLA accelerator, IMAX, in addressing the fundamental challenge of high energy consumption in LLM inference. 
We presented the first implementation of the state-of-the-art Qwen3 LLM family, along with a diverse set of quantized kernels, on IMAX using the practical llama.cpp framework. 
We then conducted a comprehensive E2E performance evaluation based on an FPGA prototype and a projected \SI{28}{\nano\meter} ASIC implementation.  
Although the proposed approach does not match the E2E latency of GPGPUs, the experimental results demonstrate its superior energy efficiency. 
The projected IMAX ASIC achieves up to a \num{44.4}$\times$ improvement in PDP and a \num{11.5}$\times$ improvement in EDP compared to the NVIDIA RTX 4090.
These results demonstrate that a general-purpose architecture can attain high energy efficiency on modern LLM inference tasks.

The analysis of the primary bottlenecks identified in this work provides clear architectural guidance for future designs. 
Our findings highlight the critical need to redesign the host-accelerator interface for data-center-scale performance and to explore co-design opportunities with emerging low-bit quantization formats. 
These promising approaches represent key directions for future research.